\chapter{Scoring de anomalía}
\label{ch:anomaly}

El ranking de caminos es la segunda mitad de lo que hace el
motor: tras matchear los caminos que satisfacen una hipótesis, el
\Scorer{} asigna a cada uno un puntaje compuesto en $[0, 1]$ y
un desglose opcional por componente para que el analista pueda
triagear. El scorer es un promedio ponderado de cinco dimensiones
(rareza de entidad, rareza de arista, concentración de vecindad,
novedad temporal y amenaza GNN), cada una computada como promedio
sobre los vértices del camino.

\section{Estado y finalización}

\begin{lstlisting}[style=rust]
pub struct AnomalyScorer {
    entity_freq:       HashMap<String, u64>,
    pair_freq:         HashMap<(String, String), u64>,
    first_seen:        HashMap<String, i64>,
    nc_cache:          HashMap<String, f64>,
    gnn_cache:         HashMap<String, f64>,
    log_max_freq:      f64,
    log_max_pair_freq: f64,
    tau_min:           i64,
    tau_max:           i64,
    weights:           ScoringWeights,
    finalized:         bool,
}
\end{lstlisting}

El scorer se puebla incrementalmente durante la ingesta
(\code{observe\_entity}, \code{observe\_edge}); son $\Ord{1}$
amortizado y no computan nada derivado. Después de la ingesta,
\code{finalize} convierte los conteos crudos en los factores de
normalización que el scoring por camino necesita.

\begin{algorithm}
\caption{\textsc{Anomaly-Finalize}}\label{alg:anomaly-finalize}
\KwIn{scorer $S$, grafo $G$}
$S.\text{log\_max\_freq} \gets \ln(\max \text{entity\_freq.values})$\;
$S.\text{log\_max\_pair\_freq} \gets \ln(\max \text{pair\_freq.values})$\;
$(\tau_{\min}, \tau_{\max}) \gets (\min, \max)$ de los valores no
cero de \code{first\_seen.values}\;
$S.\text{nc\_cache} \gets \emptyset$\;
\ForEach{$v \in G.\text{entities}$}{
    $S.\text{nc\_cache}[v] \gets $ \textsc{Compute-NC}($v, G$)\;
}
$S.\text{finalized} \gets $ verdadero\;
\end{algorithm}

\begin{complexity}
$\Ord{V + E + \sum_v d(v)} = \Ord{V + E}$ para los contadores. El
cache de entropía de Shannon cuesta un adicional
$\Ord{\sum_v d(v)} = \Ord{E}$. El espacio es $\Ord{V}$ para los
caches más lo que ya consumieron los contadores crudos.
\end{complexity}

\section{Concentración de vecindad}

\begin{algorithm}
\caption{\textsc{Compute-NC}}\label{alg:compute-nc}
\KwIn{vértice $v$, grafo $G$}
\KwOut{concentración de vecindad $\mathrm{NC}(v) \in [0, 1]$}
$\mathrm{counts} \gets \emptyset$; $\mathrm{total} \gets 0$\;
\ForEach{arista saliente $(v, u) \in G$}{
    $\mathrm{counts}[\ell_V(u)] \mathrel{+}= 1$;
    $\mathrm{total} \mathrel{+}= 1$\;
}
\ForEach{arista entrante $(w, v) \in G$}{
    $\mathrm{counts}[\ell_V(w)] \mathrel{+}= 1$;
    $\mathrm{total} \mathrel{+}= 1$\;
}
$K \gets |\mathrm{counts}|$\;
\If{$K \leq 1$}{\Return 1.0\;}
$H \gets 0$\;
\ForEach{$c \in \mathrm{counts.values}$}{
    $p \gets c / \mathrm{total}$\;
    $H \gets H - p \log_2 p$\;
}
$H_{\max} \gets \log_2 K$; $\mathrm{NC} \gets 1 - H / H_{\max}$\;
\Return clamp($\mathrm{NC}$, 0, 1)\;
\end{algorithm}

$\mathrm{NC}$ es 1 cuando la vecindad es maximamente homogénea
(una sola familia de tipo de entidad) y 0 cuando es maximamente
diversa; el clamp protege contra \emph{drift} de punto flotante
cerca de los extremos.

\begin{complexity}
$\Ord{d(v)}$ en tiempo, $\Ord{|\mathcal{T}_V|}$ en espacio (sólo
conteos por tipo de vecino; constante chica).
\end{complexity}

\section{Puntuar un camino}

Tres de las cuatro componentes no-GNN son promedios simples sobre
los vértices del camino; la componente de rareza de arista
promedia sobre las aristas del camino.

\begin{algorithm}
\caption{\textsc{Score-Path}}\label{alg:score-path}
\KwIn{camino $\pi = (v_0, \ldots, v_\ell)$, scorer $S$}
\KwOut{puntaje compuesto $c$, desglose $b$}
\If{$\pi$ vacío or $\neg S.\text{finalized}$}{\Return $(0, \mathbf{0})$\;}
\tcp{Componente 1: rareza de entidad}
$\mathrm{ER} \gets \frac{1}{\ell+1} \sum_v \left(1 - \frac{\ln(\text{entity\_freq}[v])}{S.\text{log\_max\_freq}}\right)$\;
\tcp{Componente 2: rareza de arista}
$\mathrm{EdgeR} \gets \frac{1}{\ell} \sum_{i=1}^{\ell} \left(1 - \frac{\ln(\text{pair\_freq}[v_{i-1}, v_i])}{S.\text{log\_max\_pair\_freq}}\right)$\;
\tcp{Componente 3: promedio NC}
$\mathrm{NC} \gets \frac{1}{\ell+1} \sum_v S.\text{nc\_cache}[v]$\;
\tcp{Componente 4: novedad temporal}
$\mathrm{TN} \gets \frac{1}{\ell+1} \sum_v \frac{\text{first\_seen}[v] - \tau_{\min}}{\tau_{\max} - \tau_{\min}}$\;
\tcp{Componente 5: amenaza GNN}
$\mathrm{GT} \gets \frac{1}{\ell+1} \sum_v S.\text{gnn\_cache}.\text{get}(v).\text{unwrap\_or}(0)$\;
\tcp{compuesto vía renormalización de pesos}
$W \gets \sum \text{weights}$; $c \gets \frac{1}{W} \sum_i w_i \cdot \text{componente}_i$\;
\Return (clamp($c$, 0, 1), desglose)\;
\end{algorithm}

Cada componente se clampa individualmente a $[0, 1]$ para que
una entrada de cache desactualizada (extremadamente rara, pero
posible durante un race de re-finalize) no pueda volver
negativo a un componente. La renormalización de pesos significa
que el llamador puede setearlos en cualquier escala absoluta.

\begin{complexity}
$\Ord{\ell}$ en tiempo, $\Ord{1}$ en espacio auxiliar. El estado
caro lo construyó \textsc{Anomaly-Finalize}; el scoring por
camino es barato.
\end{complexity}

\section{El \emph{seam} de trait (P2-A)}

El módulo de scoring refactorizado
(\filepath{graph\_hunter\_core/src/scoring/}) hace que cada
componente sea una \code{impl ScoreComponent} separada:

\begin{lstlisting}[style=rust]
pub trait ScoreComponent: Send + Sync {
    fn name(&self) -> &'static str;
    fn score_path(&self, ctx: &PathContext<'_>) -> f64;
    fn explain(&self, value: f64, ctx: &PathContext<'_>) -> Option<String> {
        None
    }
}
\end{lstlisting}

\code{CompositeScorer} sostiene un
\code{Vec<Box<dyn ScoreComponent>>} más pesos coincidentes y
aplica una agregación estilo \textsc{Score-Path}. Agregar un
sexto componente (bayesiano, modelo aprendido, etc.) es un
struct nuevo más un \code{vec!{}} push.

\begin{theorem}[Equivalencia v1/v2]
Para todo vector de pesos $w$ y todo camino $\pi$,
\code{AnomalyScorer::score\_path}$(w, \pi)$ y
\code{AnomalyScorer::score\_path\_v2}$(w, \pi)$ coinciden dentro
de $8 \cdot \varepsilon_{\text{f64}}$.
\end{theorem}

Esto lo verifica \code{tests/scoring\_snapshot.rs}. La
equivalencia es la precondición para eventualmente retirar el v1
hardcodeado.

\section{Estimación rápida a nivel nodo}

Durante el podado del DFS smart (Capítulo~\ref{ch:matcher})
necesitamos una estimación $\Ord{1}$ de la anomalía de un vértice
solo para que la \code{path\_anomaly\_sum} corriente pueda
rechazar ramas. Esa es \code{node\_anomaly\_estimate}:

\begin{lstlisting}[style=rust]
pub fn node_anomaly_estimate(&self, v: &str) -> f64 {
    if !self.finalized { return 0.5; }
    let er = /* Componente 1 para un solo vertice */;
    let nc = self.nc_cache.get(v).copied().unwrap_or(1.0);
    let tn = /* Componente 4 para un solo vertice */;
    let gt = self.gnn_cache.get(v).copied().unwrap_or(0.0);
    combine_with_weights(er, nc, tn, gt, &self.weights)
}
\end{lstlisting}

La rareza de arista no es parte de la estimación porque el DFS
todavía no eligió la arista siguiente cuando consulta la cota.
La estimación es entonces una cota superior suelta de la
contribución del nodo al compuesto; es admisible (nunca
sobrepoda) siempre y cuando los pesos sean no-negativos, cosa
que el motor impone al cargar la config.

\begin{inthecode}
\filepath{graph\_hunter\_core/src/anomaly.rs:148} ---
\code{finalize}.\\
\filepath{graph\_hunter\_core/src/anomaly.rs:191} ---
\code{compute\_nc}.\\
\filepath{graph\_hunter\_core/src/anomaly.rs:241} ---
\code{score\_path} (v1).\\
\filepath{graph\_hunter\_core/src/anomaly.rs:351} ---
\code{node\_anomaly\_estimate}.\\
\filepath{graph\_hunter\_core/src/scoring/} --- v2 basado en trait.\\
\filepath{graph\_hunter\_core/tests/scoring\_snapshot.rs} ---
fijación de equivalencia v1/v2.
\end{inthecode}
